% !TEX root = ../../main.tex
\chapter*{前两章概念复查表}
\phantomsection
\addcontentsline{toc}{chapter}{前两章概念复查表}
\markboth{前两章概念复查表}{前两章概念复查表}

按正文顺序列出第一章、第二章已经引入的概念。复查时先尝试完整说出定义及其条件，再核对要点；点击概念名称或页码可跳转到定义位置。方框可用于打印后勾选。同一处同时定义的相关概念共用跳转位置。

\begingroup
\small
\setlength{\tabcolsep}{3pt}
\renewcommand{\arraystretch}{1.22}
\begin{longtable}{@{}p{0.03\linewidth}>{\raggedright\arraybackslash}p{0.255\linewidth}>{\raggedright\arraybackslash}p{0.55\linewidth}>{\raggedright\arraybackslash}p{0.10\linewidth}@{}}
\hline
 & \textbf{概念} & \textbf{复查要点} & \textbf{定义位置}\\
\hline
\endfirsthead
\multicolumn{4}{r}{\small 前两章概念复查表（续）}\\
\hline
 & \textbf{概念} & \textbf{复查要点} & \textbf{定义位置}\\
\hline
\endhead
\hline
\multicolumn{4}{r}{\small 续下页}\\
\endfoot
\hline
\endlastfoot
\multicolumn{4}{@{}l@{}}{\rule{0pt}{2.4em}\large\textbf{第一章：代数基本概念}}\\*
\multicolumn{4}{@{}l@{}}{\rule{0pt}{1.8em}\textbf{§1.1\quad 代数运算}}\\*
$\square$ & \hyperref[p2:concept:c1-l19]{代数运算（二元运算）} & $A\times A\to A$，结果唯一且仍在$A$中。 & \hyperref[p2:concept:c1-l19]{第\pageref*{p2:concept:c1-l19}页}\\
$\square$ & \hyperref[p2:concept:c1-l34]{结合律} & $(a*b)*c=a*(b*c)$。 & \hyperref[p2:concept:c1-l34]{第\pageref*{p2:concept:c1-l34}页}\\
$\square$ & \hyperref[p2:concept:c1-l37]{交换律} & $a*b=b*a$。 & \hyperref[p2:concept:c1-l37]{第\pageref*{p2:concept:c1-l37}页}\\
$\square$ & \hyperref[p2:concept:c1-l40]{单位元} & $e*a=a*e=a$，对所有$a$成立。 & \hyperref[p2:concept:c1-l40]{第\pageref*{p2:concept:c1-l40}页}\\
$\square$ & \hyperref[p2:concept:c1-l42]{逆元} & $a*b=b*a=e$，以单位元已经存在为前提。 & \hyperref[p2:concept:c1-l42]{第\pageref*{p2:concept:c1-l42}页}\\
\multicolumn{4}{@{}l@{}}{\rule{0pt}{1.8em}\textbf{§1.2\quad 群的定义和简单性质}}\\*
$\square$ & \hyperref[p2:concept:c1-l52]{群} & 运算结合，存在左单位元，每个元素有左逆元。 & \hyperref[p2:concept:c1-l52]{第\pageref*{p2:concept:c1-l52}页}\\
$\square$ & \hyperref[p2:concept:c1-l55]{左单位元} & $ea=a$，对所有$a\in G$成立。 & \hyperref[p2:concept:c1-l55]{第\pageref*{p2:concept:c1-l55}页}\\
$\square$ & \hyperref[p2:concept:c1-l56]{左逆元} & $ba=e$。 & \hyperref[p2:concept:c1-l56]{第\pageref*{p2:concept:c1-l56}页}\\
$\square$ & \hyperref[p2:concept:c1-l60]{交换群（阿贝尔群）} & 群中任意两个元素交换。 & \hyperref[p2:concept:c1-l60]{第\pageref*{p2:concept:c1-l60}页}\\
$\square$ & \hyperref[p2:concept:c1-l105]{元素的整数次幂} & $a^0=e$，$a^{-n}=(a^{-1})^n$。 & \hyperref[p2:concept:c1-l105]{第\pageref*{p2:concept:c1-l105}页}\\
$\square$ & \hyperref[p2:concept:c1-l106]{加法记号下的整数倍} & $na$对应乘法记号下的$a^n$。 & \hyperref[p2:concept:c1-l106]{第\pageref*{p2:concept:c1-l106}页}\\
$\square$ & \hyperref[p2:concept:c1-l149]{群的阶} & $\lvert G\rvert$为元素个数。 & \hyperref[p2:concept:c1-l149]{第\pageref*{p2:concept:c1-l149}页}\\
$\square$ & \hyperref[p2:concept:c1-l149]{有限群、无限群} & 按元素个数是否有限区分。 & \hyperref[p2:concept:c1-l149]{第\pageref*{p2:concept:c1-l149}页}\\
\multicolumn{4}{@{}l@{}}{\rule{0pt}{1.8em}\textbf{§1.3\quad 群的例子}}\\*
$\square$ & \hyperref[p2:concept:c1-l156]{图形的对称群} & 保持图形整体不变的等距变换，在复合下成群。 & \hyperref[p2:concept:c1-l156]{第\pageref*{p2:concept:c1-l156}页}\\
$\square$ & \hyperref[p2:concept:c1-l164]{二面体群} & $D_n$为正$n$边形的对称群，阶为$2n$。 & \hyperref[p2:concept:c1-l164]{第\pageref*{p2:concept:c1-l164}页}\\
$\square$ & \hyperref[p2:concept:c1-l192]{全变换群} & $S(M)$为集合$M$的全部双射，运算为复合。 & \hyperref[p2:concept:c1-l192]{第\pageref*{p2:concept:c1-l192}页}\\
$\square$ & \hyperref[p2:concept:c1-l193]{$n$元对称群} & $S_n=S(\{1,\ldots,n\})$。 & \hyperref[p2:concept:c1-l193]{第\pageref*{p2:concept:c1-l193}页}\\
$\square$ & \hyperref[p2:concept:c1-l194]{置换} & 对称群的元素，复合约定为右边先作用。 & \hyperref[p2:concept:c1-l194]{第\pageref*{p2:concept:c1-l194}页}\\
$\square$ & \hyperref[p2:concept:c1-l209]{轮换} & 依次循环指定元素，其余元素保持不变。 & \hyperref[p2:concept:c1-l209]{第\pageref*{p2:concept:c1-l209}页}\\
$\square$ & \hyperref[p2:concept:c1-l210]{对换} & 长度为$2$的轮换。 & \hyperref[p2:concept:c1-l210]{第\pageref*{p2:concept:c1-l210}页}\\
$\square$ & \hyperref[p2:concept:c1-l210]{不相交轮换} & 各轮换涉及的元素集合互不相交。 & \hyperref[p2:concept:c1-l210]{第\pageref*{p2:concept:c1-l210}页}\\
\multicolumn{4}{@{}l@{}}{\rule{0pt}{1.8em}\textbf{§1.4\quad 子群、陪集}}\\*
$\square$ & \hyperref[p2:concept:c1-l247]{子群} & 非空子集在原运算下成群，记为$H\leq G$。 & \hyperref[p2:concept:c1-l247]{第\pageref*{p2:concept:c1-l247}页}\\
$\square$ & \hyperref[p2:concept:c1-l248]{变换群} & 全变换群的子群。 & \hyperref[p2:concept:c1-l248]{第\pageref*{p2:concept:c1-l248}页}\\
$\square$ & \hyperref[p2:concept:c1-l252]{平凡子群} & $\{e\}$和$G$。 & \hyperref[p2:concept:c1-l252]{第\pageref*{p2:concept:c1-l252}页}\\
$\square$ & \hyperref[p2:concept:c1-l253]{非平凡子群} & 同时不同于$\{e\}$和$G$的子群。 & \hyperref[p2:concept:c1-l253]{第\pageref*{p2:concept:c1-l253}页}\\
$\square$ & \hyperref[p2:concept:c1-l284]{置换的符号} & $\operatorname{sgn}(\sigma)$由差积变换后的符号确定。 & \hyperref[p2:concept:c1-l284]{第\pageref*{p2:concept:c1-l284}页}\\
$\square$ & \hyperref[p2:concept:c1-l289]{偶置换、奇置换} & 符号分别为$1$、$-1$。 & \hyperref[p2:concept:c1-l289]{第\pageref*{p2:concept:c1-l289}页}\\
$\square$ & \hyperref[p2:concept:c1-l290]{交错群} & $A_n$为全体偶置换组成的群。 & \hyperref[p2:concept:c1-l290]{第\pageref*{p2:concept:c1-l290}页}\\
$\square$ & \hyperref[p2:concept:c1-l311]{左陪集、右陪集} & $aH=\{ah:h\in H\}$，$Ha=\{ha:h\in H\}$。 & \hyperref[p2:concept:c1-l311]{第\pageref*{p2:concept:c1-l311}页}\\
$\square$ & \hyperref[p2:concept:c1-l312]{陪集代表} & 表示陪集$aH$或$Ha$的元素$a$。 & \hyperref[p2:concept:c1-l312]{第\pageref*{p2:concept:c1-l312}页}\\
$\square$ & \hyperref[p2:concept:c1-l312]{子群的指数} & $[G:H]$为左陪集的个数。 & \hyperref[p2:concept:c1-l312]{第\pageref*{p2:concept:c1-l312}页}\\
$\square$ & \hyperref[p2:concept:c1-l353]{由单个元素生成的子群} & $\langle a\rangle=\{a^m:m\in\Z\}$。 & \hyperref[p2:concept:c1-l353]{第\pageref*{p2:concept:c1-l353}页}\\
$\square$ & \hyperref[p2:concept:c1-l354]{元素的阶} & $\operatorname{ord}(a)=\lvert\langle a\rangle\rvert$。 & \hyperref[p2:concept:c1-l354]{第\pageref*{p2:concept:c1-l354}页}\\
\multicolumn{4}{@{}l@{}}{\rule{0pt}{1.8em}\textbf{§1.5\quad 群的同构}}\\*
$\square$ & \hyperref[p2:concept:c1-l388]{群同构、同构的群} & 保持乘法的双射，记为$G\cong G'$。 & \hyperref[p2:concept:c1-l388]{第\pageref*{p2:concept:c1-l388}页}\\
$\square$ & \hyperref[p2:concept:c1-l468]{左正则表示、右正则表示} & 分别由$L_a(x)=ax$、$R_a(x)=xa^{-1}$得到。 & \hyperref[p2:concept:c1-l468]{第\pageref*{p2:concept:c1-l468}页}\\
\multicolumn{4}{@{}l@{}}{\rule{0pt}{1.8em}\textbf{§1.6\quad 同态、正规子群}}\\*
$\square$ & \hyperref[p2:concept:c1-l482]{群同态} & 映射满足$\varphi(ab)=\varphi(a)\varphi(b)$。 & \hyperref[p2:concept:c1-l482]{第\pageref*{p2:concept:c1-l482}页}\\
$\square$ & \hyperref[p2:concept:c1-l483]{群同态的像} & $\varphi(G)=\{\varphi(a):a\in G\}$。 & \hyperref[p2:concept:c1-l483]{第\pageref*{p2:concept:c1-l483}页}\\
$\square$ & \hyperref[p2:concept:c1-l483]{群同态的核} & $\ker\varphi=\{a:\varphi(a)=e'\}$。 & \hyperref[p2:concept:c1-l483]{第\pageref*{p2:concept:c1-l483}页}\\
$\square$ & \hyperref[p2:concept:c1-l487]{群的嵌入} & 单射的群同态。 & \hyperref[p2:concept:c1-l487]{第\pageref*{p2:concept:c1-l487}页}\\
$\square$ & \hyperref[p2:concept:c1-l487]{群的满同态} & 满射的群同态。 & \hyperref[p2:concept:c1-l487]{第\pageref*{p2:concept:c1-l487}页}\\
$\square$ & \hyperref[p2:concept:c1-l488]{元素的完全反像} & $\varphi^{-1}(\{y\})$是一个集合，不要求存在逆映射。 & \hyperref[p2:concept:c1-l488]{第\pageref*{p2:concept:c1-l488}页}\\
$\square$ & \hyperref[p2:concept:c1-l574]{正规子群} & $aH=Ha$对所有$a\in G$成立，记为$H\trianglelefteq G$。 & \hyperref[p2:concept:c1-l574]{第\pageref*{p2:concept:c1-l574}页}\\
\multicolumn{4}{@{}l@{}}{\rule{0pt}{1.8em}\textbf{§1.7\quad 商群}}\\*
$\square$ & \hyperref[p2:concept:c1-l613]{子集的乘积} & $AB=\{ab:a\in A,b\in B\}$。 & \hyperref[p2:concept:c1-l613]{第\pageref*{p2:concept:c1-l613}页}\\
$\square$ & \hyperref[p2:concept:c1-l613]{子集的逆集} & $A^{-1}=\{a^{-1}:a\in A\}$。 & \hyperref[p2:concept:c1-l613]{第\pageref*{p2:concept:c1-l613}页}\\
$\square$ & \hyperref[p2:concept:c1-l681]{商群} & 正规子群$H$的陪集全体，乘法为$(aH)(bH)=abH$。 & \hyperref[p2:concept:c1-l681]{第\pageref*{p2:concept:c1-l681}页}\\
$\square$ & \hyperref[p2:concept:c1-l727]{群的自然同态} & $\pi:G\to G/H$，$\pi(a)=aH$。 & \hyperref[p2:concept:c1-l727]{第\pageref*{p2:concept:c1-l727}页}\\
$\square$ & \hyperref[p2:concept:c1-l762]{交换图} & 同起点、同终点的路径给出相同的复合映射。 & \hyperref[p2:concept:c1-l762]{第\pageref*{p2:concept:c1-l762}页}\\
\multicolumn{4}{@{}l@{}}{\rule{0pt}{1.8em}\textbf{§1.8\quad 环、子环}}\\*
$\square$ & \hyperref[p2:concept:c1-l781]{环} & 加法为交换群，乘法结合，并满足两侧分配律。 & \hyperref[p2:concept:c1-l781]{第\pageref*{p2:concept:c1-l781}页}\\
$\square$ & \hyperref[p2:concept:c1-l785]{环的零元、加法逆元} & 分别记为$0$与$-a$。 & \hyperref[p2:concept:c1-l785]{第\pageref*{p2:concept:c1-l785}页}\\
$\square$ & \hyperref[p2:concept:c1-l830]{环中的整数倍、正整数幂} & 整数倍来自加法群，正整数幂来自连乘。 & \hyperref[p2:concept:c1-l830]{第\pageref*{p2:concept:c1-l830}页}\\
$\square$ & \hyperref[p2:concept:c1-l887]{子环} & 非空子集在原有两种运算下成环。 & \hyperref[p2:concept:c1-l887]{第\pageref*{p2:concept:c1-l887}页}\\
$\square$ & \hyperref[p2:concept:c1-l923]{环同构} & 同时保持加法和乘法的双射。 & \hyperref[p2:concept:c1-l923]{第\pageref*{p2:concept:c1-l923}页}\\
\multicolumn{4}{@{}l@{}}{\rule{0pt}{1.8em}\textbf{§1.9\quad 各种特殊类型的环}}\\*
$\square$ & \hyperref[p2:concept:c1-l954]{幺环、乘法单位元} & 存在乘法单位元$1$的环。 & \hyperref[p2:concept:c1-l954]{第\pageref*{p2:concept:c1-l954}页}\\
$\square$ & \hyperref[p2:concept:c1-l956]{单位（可逆元）} & 有双边乘法逆元的元素，全体记为$R^\times$。 & \hyperref[p2:concept:c1-l956]{第\pageref*{p2:concept:c1-l956}页}\\
$\square$ & \hyperref[p2:concept:c1-l958]{交换环} & 乘法满足交换律的环。 & \hyperref[p2:concept:c1-l958]{第\pageref*{p2:concept:c1-l958}页}\\
$\square$ & \hyperref[p2:concept:c1-l963]{单位群} & $R^\times$在乘法下组成的群。 & \hyperref[p2:concept:c1-l963]{第\pageref*{p2:concept:c1-l963}页}\\
$\square$ & \hyperref[p2:concept:c1-l988]{左零因子} & 非零$a$使$ab=0$对某个非零$b$成立。 & \hyperref[p2:concept:c1-l988]{第\pageref*{p2:concept:c1-l988}页}\\
$\square$ & \hyperref[p2:concept:c1-l989]{右零因子} & 非零$a$使$ba=0$对某个非零$b$成立。 & \hyperref[p2:concept:c1-l989]{第\pageref*{p2:concept:c1-l989}页}\\
$\square$ & \hyperref[p2:concept:c1-l989]{零因子} & 左零因子与右零因子的统称。 & \hyperref[p2:concept:c1-l989]{第\pageref*{p2:concept:c1-l989}页}\\
$\square$ & \hyperref[p2:concept:c1-l990]{整环} & 无零因子的交换幺环，且$1\ne0$。 & \hyperref[p2:concept:c1-l990]{第\pageref*{p2:concept:c1-l990}页}\\
$\square$ & \hyperref[p2:concept:c1-l1035]{域} & 交换幺环，$1\ne0$，每个非零元素可逆。 & \hyperref[p2:concept:c1-l1035]{第\pageref*{p2:concept:c1-l1035}页}\\
$\square$ & \hyperref[p2:concept:c1-l1037]{体（除环）} & 保留域的其他条件，去掉乘法交换的要求。 & \hyperref[p2:concept:c1-l1037]{第\pageref*{p2:concept:c1-l1037}页}\\
$\square$ & \hyperref[p2:concept:c1-l1039]{子域} & 域的子环，而且其本身也是域。 & \hyperref[p2:concept:c1-l1039]{第\pageref*{p2:concept:c1-l1039}页}\\
$\square$ & \hyperref[p2:concept:c1-l1092]{四元数体} & 用$X(\alpha,\beta)$矩阵模型定义的非交换体$\mathbb H$。 & \hyperref[p2:concept:c1-l1092]{第\pageref*{p2:concept:c1-l1092}页}\\
\multicolumn{4}{@{}l@{}}{\rule{0pt}{1.8em}\textbf{§1.10\quad 环的同态、理想}}\\*
$\square$ & \hyperref[p2:concept:c1-l1152]{环同态} & 同时保持加法、乘法，不预先要求保持$1$。 & \hyperref[p2:concept:c1-l1152]{第\pageref*{p2:concept:c1-l1152}页}\\
$\square$ & \hyperref[p2:concept:c1-l1156]{环同态的像} & $\varphi(R)$为同态的全部取值组成的子环。 & \hyperref[p2:concept:c1-l1156]{第\pageref*{p2:concept:c1-l1156}页}\\
$\square$ & \hyperref[p2:concept:c1-l1156]{环同态的核} & $\ker\varphi=\{a:\varphi(a)=0\}$。 & \hyperref[p2:concept:c1-l1156]{第\pageref*{p2:concept:c1-l1156}页}\\
$\square$ & \hyperref[p2:concept:c1-l1157]{零同态} & 像为$\{0\}$的环同态。 & \hyperref[p2:concept:c1-l1157]{第\pageref*{p2:concept:c1-l1157}页}\\
$\square$ & \hyperref[p2:concept:c1-l1157]{环的满同态} & 像为整个目标环的环同态。 & \hyperref[p2:concept:c1-l1157]{第\pageref*{p2:concept:c1-l1157}页}\\
$\square$ & \hyperref[p2:concept:c1-l1201]{双边理想（理想）} & 加法子群，且被母环元素左右吸收。 & \hyperref[p2:concept:c1-l1201]{第\pageref*{p2:concept:c1-l1201}页}\\
$\square$ & \hyperref[p2:concept:c1-l1202]{左理想、右理想} & 分别只要求$ra\in I$、$ar\in I$。 & \hyperref[p2:concept:c1-l1202]{第\pageref*{p2:concept:c1-l1202}页}\\
$\square$ & \hyperref[p2:concept:c1-l1203]{平凡理想} & $\{0\}$与$R$。 & \hyperref[p2:concept:c1-l1203]{第\pageref*{p2:concept:c1-l1203}页}\\
\multicolumn{4}{@{}l@{}}{\rule{0pt}{1.8em}\textbf{§1.11\quad 商环}}\\*
$\square$ & \hyperref[p2:concept:c1-l1252]{商环} & 双边理想的加法陪集全体，按代表元的和与积定义运算。 & \hyperref[p2:concept:c1-l1252]{第\pageref*{p2:concept:c1-l1252}页}\\
$\square$ & \hyperref[p2:concept:c1-l1261]{零环} & 只有零元的环，这里由$R/R$得到。 & \hyperref[p2:concept:c1-l1261]{第\pageref*{p2:concept:c1-l1261}页}\\
$\square$ & \hyperref[p2:concept:c1-l1328]{环的自然同态} & $\pi:R\to R/I$，$\pi(a)=a+I$。 & \hyperref[p2:concept:c1-l1328]{第\pageref*{p2:concept:c1-l1328}页}\\
$\square$ & \hyperref[p2:concept:c1-l1356]{欧拉函数} & $\varphi(n)$为模$n$单位群的阶，即与$n$互素的剩余类数。 & \hyperref[p2:concept:c1-l1356]{第\pageref*{p2:concept:c1-l1356}页}\\
\multicolumn{4}{@{}l@{}}{\rule{0pt}{1.8em}\textbf{§1.12\quad 特征}}\\*
$\square$ & \hyperref[p2:concept:c1-l1442]{域的特征} & $1$的有限加法阶，若其加法阶无限则约定为$0$。 & \hyperref[p2:concept:c1-l1442]{第\pageref*{p2:concept:c1-l1442}页}\\
$\square$ & \hyperref[p2:concept:c1-l1517]{素域} & 每个域中唯一的最小子域。 & \hyperref[p2:concept:c1-l1517]{第\pageref*{p2:concept:c1-l1517}页}\\
\multicolumn{4}{@{}l@{}}{\rule{0pt}{2.4em}\large\textbf{第二章：群}}\\*
\multicolumn{4}{@{}l@{}}{\rule{0pt}{1.8em}\textbf{§2.1\quad 群的同态定理}}\\*
$\square$ & \hyperref[p2:concept:c2-l12]{子群的完全反像} & $\pi^{-1}(\overline H)=\{g:\pi(g)\in\overline H\}$。 & \hyperref[p2:concept:c2-l12]{第\pageref*{p2:concept:c2-l12}页}\\
$\square$ & \hyperref[p2:concept:c2-l142]{由集合生成的子群} & $\langle S\rangle$为包含$S$的全部子群之交。 & \hyperref[p2:concept:c2-l142]{第\pageref*{p2:concept:c2-l142}页}\\
$\square$ & \hyperref[p2:concept:c2-l146]{生成元（生成元组）} & 满足$\langle S\rangle=G$的集合$S$。 & \hyperref[p2:concept:c2-l146]{第\pageref*{p2:concept:c2-l146}页}\\
$\square$ & \hyperref[p2:concept:c2-l146]{有限生成群} & 存在有限生成元组的群。 & \hyperref[p2:concept:c2-l146]{第\pageref*{p2:concept:c2-l146}页}\\
$\square$ & \hyperref[p2:concept:c2-l173]{换位子} & 约定$[a,b]=a^{-1}b^{-1}ab$。 & \hyperref[p2:concept:c2-l173]{第\pageref*{p2:concept:c2-l173}页}\\
$\square$ & \hyperref[p2:concept:c2-l173]{换位子群} & 全体换位子生成的子群$G^{(1)}=[G,G]$。 & \hyperref[p2:concept:c2-l173]{第\pageref*{p2:concept:c2-l173}页}\\
\multicolumn{4}{@{}l@{}}{\rule{0pt}{1.8em}\textbf{§2.2\quad 循环群}}\\*
$\square$ & \hyperref[p2:concept:c2-l204]{循环群} & 由一个元素生成的群$G=\langle g\rangle$。 & \hyperref[p2:concept:c2-l204]{第\pageref*{p2:concept:c2-l204}页}\\
\multicolumn{4}{@{}l@{}}{\rule{0pt}{1.8em}\textbf{§2.3\quad 单群与交错群的单性}}\\*
$\square$ & \hyperref[p2:concept:c2-l332]{单群} & 非平凡群，只有$\{e\}$和$G$两个正规子群。 & \hyperref[p2:concept:c2-l332]{第\pageref*{p2:concept:c2-l332}页}\\
$\square$ & \hyperref[p2:concept:c2-l372]{轮换类型} & 计入长度$1$的轮换，记录全部轮换长度，保留重复且不计次序。 & \hyperref[p2:concept:c2-l372]{第\pageref*{p2:concept:c2-l372}页}\\
$\square$ & \hyperref[p2:concept:c2-l405]{置换的不动点、被移动的点} & 分别满足$\sigma(i)=i$、$\sigma(i)\ne i$。 & \hyperref[p2:concept:c2-l405]{第\pageref*{p2:concept:c2-l405}页}\\
\multicolumn{4}{@{}l@{}}{\rule{0pt}{1.8em}\textbf{§2.4\quad 可解群}}\\*
$\square$ & \hyperref[p2:concept:c2-l504]{导出子群与导出列} & $G^{(0)}=G$，$G^{(k)}=(G^{(k-1)})^{(1)}$。 & \hyperref[p2:concept:c2-l504]{第\pageref*{p2:concept:c2-l504}页}\\
$\square$ & \hyperref[p2:concept:c2-l522]{可解群} & 存在正整数$k$使$G^{(k)}=\{e\}$。 & \hyperref[p2:concept:c2-l522]{第\pageref*{p2:concept:c2-l522}页}\\
$\square$ & \hyperref[p2:concept:c2-l617]{极大正规子群} & 真正规子群，且与$G$之间无其他$G$的正规子群。 & \hyperref[p2:concept:c2-l617]{第\pageref*{p2:concept:c2-l617}页}\\
\multicolumn{4}{@{}l@{}}{\rule{0pt}{1.8em}\textbf{§2.5\quad 群的自同构群}}\\*
$\square$ & \hyperref[p2:concept:c2-l718]{自同构} & 群到自身的同构。 & \hyperref[p2:concept:c2-l718]{第\pageref*{p2:concept:c2-l718}页}\\
$\square$ & \hyperref[p2:concept:c2-l722]{自同构群} & $\Aut(G)$，乘法为映射复合。 & \hyperref[p2:concept:c2-l722]{第\pageref*{p2:concept:c2-l722}页}\\
$\square$ & \hyperref[p2:concept:c2-l790]{内自同构} & 由群内元素$a$引起的共轭自同构$\sigma_a$。 & \hyperref[p2:concept:c2-l790]{第\pageref*{p2:concept:c2-l790}页}\\
$\square$ & \hyperref[p2:concept:c2-l799]{内自同构群} & $\operatorname{Inn}(G)$，由全部内自同构组成。 & \hyperref[p2:concept:c2-l799]{第\pageref*{p2:concept:c2-l799}页}\\
$\square$ & \hyperref[p2:concept:c2-l808]{群的中心} & $Z(G)=\{a:ax=xa\text{ 对所有 }x\in G\}$。 & \hyperref[p2:concept:c2-l808]{第\pageref*{p2:concept:c2-l808}页}\\
$\square$ & \hyperref[p2:concept:c2-l845]{外自同构群} & $\operatorname{Out}(G)=\Aut(G)/\operatorname{Inn}(G)$。 & \hyperref[p2:concept:c2-l845]{第\pageref*{p2:concept:c2-l845}页}\\
$\square$ & \hyperref[p2:concept:c2-l885]{完全群} & 中心平凡，且每个自同构都是内自同构。 & \hyperref[p2:concept:c2-l885]{第\pageref*{p2:concept:c2-l885}页}\\
$\square$ & \hyperref[p2:concept:c2-l891]{特征子群} & $\alpha(H)=H$对每个$\alpha\in\Aut(G)$成立。 & \hyperref[p2:concept:c2-l891]{第\pageref*{p2:concept:c2-l891}页}\\
\multicolumn{4}{@{}l@{}}{\rule{0pt}{1.8em}\textbf{§2.6\quad 群在集合上的作用}}\\*
$\square$ & \hyperref[p2:concept:c2-l924]{群的左作用} & $e(x)=x$且$g_1(g_2(x))=(g_1g_2)(x)$。 & \hyperref[p2:concept:c2-l924]{第\pageref*{p2:concept:c2-l924}页}\\
$\square$ & \hyperref[p2:concept:c2-l963]{左平移作用} & 在$G$自身上定义$g(x)=gx$。 & \hyperref[p2:concept:c2-l963]{第\pageref*{p2:concept:c2-l963}页}\\
$\square$ & \hyperref[p2:concept:c2-l971]{共轭作用} & 在$G$自身上定义$g(x)=gxg^{-1}$。 & \hyperref[p2:concept:c2-l971]{第\pageref*{p2:concept:c2-l971}页}\\
$\square$ & \hyperref[p2:concept:c2-l978]{共轭元素、共轭子群} & 分别形如$gxg^{-1}$、$gHg^{-1}$。 & \hyperref[p2:concept:c2-l978]{第\pageref*{p2:concept:c2-l978}页}\\
$\square$ & \hyperref[p2:concept:c2-l986]{齐性空间} & 全部左陪集组成的集合$G/H$，不要求$H$正规。 & \hyperref[p2:concept:c2-l986]{第\pageref*{p2:concept:c2-l986}页}\\
$\square$ & \hyperref[p2:concept:c2-l1001]{忠实（如实）作用} & 对应的同态$G\to S(X)$为单射。 & \hyperref[p2:concept:c2-l1001]{第\pageref*{p2:concept:c2-l1001}页}\\
$\square$ & \hyperref[p2:concept:c2-l1004]{等价的作用} & 存在双射$\varphi$满足$\varphi(g(x))=g(\varphi(x))$。 & \hyperref[p2:concept:c2-l1004]{第\pageref*{p2:concept:c2-l1004}页}\\
$\square$ & \hyperref[p2:concept:c2-l1031]{轨道} & $O_x=\{g(x):g\in G\}$。 & \hyperref[p2:concept:c2-l1031]{第\pageref*{p2:concept:c2-l1031}页}\\
$\square$ & \hyperref[p2:concept:c2-l1035]{作用的不动点及全体不动点} & $g(x)=x$对所有$g$成立，全体记为$X^G$。 & \hyperref[p2:concept:c2-l1035]{第\pageref*{p2:concept:c2-l1035}页}\\
$\square$ & \hyperref[p2:concept:c2-l1035]{传递作用} & 整个$X$只有一个轨道。 & \hyperref[p2:concept:c2-l1035]{第\pageref*{p2:concept:c2-l1035}页}\\
$\square$ & \hyperref[p2:concept:c2-l1046]{稳定子群} & $H_x=\{g\in G:g(x)=x\}$。 & \hyperref[p2:concept:c2-l1046]{第\pageref*{p2:concept:c2-l1046}页}\\
$\square$ & \hyperref[p2:concept:c2-l1111]{有限$p$群} & 群阶为$p^k$，其中$p$为素数、$k\ge1$。 & \hyperref[p2:concept:c2-l1111]{第\pageref*{p2:concept:c2-l1111}页}\\
$\square$ & \hyperref[p2:concept:c2-l1151]{共轭类} & $C(x)=\{gxg^{-1}:g\in G\}$。 & \hyperref[p2:concept:c2-l1151]{第\pageref*{p2:concept:c2-l1151}页}\\
$\square$ & \hyperref[p2:concept:c2-l1162]{中心化子} & $C_G(x)$为与$x$交换的全部群元素组成的子群。 & \hyperref[p2:concept:c2-l1162]{第\pageref*{p2:concept:c2-l1162}页}\\
$\square$ & \hyperref[p2:concept:c2-l1208]{正规化子} & $N_G(H)=\{g:gHg^{-1}=H\}$。 & \hyperref[p2:concept:c2-l1208]{第\pageref*{p2:concept:c2-l1208}页}\\
\multicolumn{4}{@{}l@{}}{\rule{0pt}{1.8em}\textbf{§2.7\quad 西罗定理}}\\*
$\square$ & \hyperref[p2:concept:c2-l1287]{西罗$p$子群} & 若$\lvert G\rvert=p^lm$且$p\nmid m$，则为$p^l$阶子群。 & \hyperref[p2:concept:c2-l1287]{第\pageref*{p2:concept:c2-l1287}页}\\
\multicolumn{4}{@{}l@{}}{\rule{0pt}{1.8em}\textbf{§2.8\quad 群的直和}}\\*
$\square$ & \hyperref[p2:concept:c2-l1501]{外直和} & $G_1\oplus\cdots\oplus G_s$按分量定义乘法。 & \hyperref[p2:concept:c2-l1501]{第\pageref*{p2:concept:c2-l1501}页}\\
$\square$ & \hyperref[p2:concept:c2-l1505]{有限个群的直积} & 本章有限个因子的直积与外直和是同一构造。 & \hyperref[p2:concept:c2-l1505]{第\pageref*{p2:concept:c2-l1505}页}\\
$\square$ & \hyperref[p2:concept:c2-l1542]{内直和} & 正规子群的乘积等于$G$，且每个元素的分量表示唯一。 & \hyperref[p2:concept:c2-l1542]{第\pageref*{p2:concept:c2-l1542}页}\\
$\square$ & \hyperref[p2:concept:c2-l1600]{不可分解群} & 非平凡群，不能分解为两个非平凡群的直和。 & \hyperref[p2:concept:c2-l1600]{第\pageref*{p2:concept:c2-l1600}页}\\
\end{longtable}
\endgroup
